\section{Conclusion}\label{sec:conclusion}

We asked how an LLM represents what different people find dramatic, and we find a split answer that cleanly separates the structure of the construct from the structure of disagreement. On the construct side, the model represents drama as a low-dimensional \emph{space} of features rather than a single axis: drama itself decodes perfectly (AUROC $1.000$), but melodrama is not just more drama---the dramatic$\to$melodramatic step rotates to $\sim$83 degrees from the drama axis in the deepest layers (\Tabref{tab:e1}), dramatic and melodramatic texts sit at the same drama-axis value while splitting along an orthogonal \emph{excess} direction (\figref{fig:excess}), and an 8-subconcept subspace has participation ratio rising to $\sim$5.6 ($\sim$5.9 in \qwensmall), far above a single axis (\figref{fig:geometry}). ``Different kinds of drama'' are therefore genuine and geometric. On the disagreement side, the picture is the opposite: per-reader variation is overwhelmingly a \emph{threshold} on a shared drama axis. Adding a per-reader threshold lifts held-out AUROC from $0.510$ to $0.619$ (transferred subspace) and from $0.660$ to $0.700$ (in-domain), while a shared full subspace barely helps beyond one axis ($\sim+0.001$) and letting each reader pick their own \emph{direction} adds exactly the shuffled-reader null ($\Delta$AUROC $+0.0076$ real vs.\ $+0.0076$ null, $z\!\approx\!1.0$, n.s.; \Tabref{tab:e3indomain}, \figref{fig:e3b}). Six prompted personas leave the drama axis essentially fixed (pairwise cosine $\geq 0.974$, mean $0.99$; \figref{fig:persona}). The strong multi-dimensional reader hypothesis is not supported, for either natural disagreement or prompted personas. {\bf The key takeaway: drama's multidimensionality lives in the shared content, while inter-reader differences are mostly about where the threshold sits.}

Several directions would give the reader-direction hypothesis its strongest possible chance and pin down the excess axis causally. First, our \goemotions\ labels are binary high-arousal indicators, so our verdict is ``no \emph{detectable} reader-specific direction''; a purpose-built, graded, multi-reader \emph{drama} rating set could expose reader-specific structure that a binary indicator cannot resolve. Second, causal steering along the drama and excess axes \citep{wollschlager2025geometry} would confirm that the two directions are independent under intervention, not merely in observed geometry, and SAE or feature-circuit localization \citep{marks2024sfc} could identify the components that implement the excess direction. Third, our results hold across \qwensmall\ and \qwenbig\ and across layers and subspace dimensions, but cross-family replication on Llama and Gemma is needed to rule out a Qwen-specific geometry. Finally, prompted personas did not rotate the drama axis; a natural next test is whether \emph{fine-tuned} reader embeddings---rather than system prompts---can rotate the axis where personas could not, which would be the most direct route from a shared-threshold account back toward genuinely reader-specific directions.
